\section[SNC 因子と connection holonomy による有限相対 compactification]{SNC 因子と connection holonomy による\\有限相対 compactification}
\label{snc:section}
% 2026-11-21 / v1.2: 固定 connection の全 leaf-period 条件と最小 holonomy 自由度．

直線配置に対する構成の本質は，個々の成分が ample であることではない．
必要なのは，Kähler class が divisor 全体の正の重み付き和で表されることである．
有理曲線の場合の graph periods に加え，正の genus を持つ成分上の periods は
normal connection の holonomy を選ぶことで解消できる．
本節では，任意の genus，負の自己交点数，同じ二成分の複数交点，
および二価の成分を含む形でこの構成を証明する．

\subsection{仮定，marking，結論}

$M$ を連結な滑らかな射影複素曲面，$D=\bigcup_{a=1}^{r}D_a$ を
reduced simple normal crossings divisor とする．
各 $D_a$ は genus $g_a$ の滑らかな曲線とし，三重交点はない．
異なる二成分が複数の点で交わることを許す．
各交点を個別にラベル付けし，その集合を $E$ とする．
$e=|E|$，$n_a=|D_a\cap\bigcup_{b\ne a}D_b|$ と置く．
$n_a$ は隣接する成分の数ではなく，\emph{交点の数}である．
以下では
\begin{equation}\label{snc:valence}
 n_a\geq2\qquad(1\leq a\leq r)
\end{equation}
を仮定する．$G$ は頂点が $D_a$，辺が個々の交点である dual multigraph とする．

Kähler form $\omega$ と正の実数 $w_a$ を，
\begin{equation}\label{snc:positive-class}
 [\omega]=2\pi\sum_{a=1}^{r}w_a[D_a],\qquad w_a>0
\end{equation}
となるように固定する．重みは存在だけでなくデータの一部として固定する．
normal-circle の周期は $2\pi$，Hamiltonian convention は
$\iota_V\omega=-\dd\mu$，normal moment は $\mu\geq0$ とする．
angular connection の curvature convention は
\begin{equation}\label{snc:curvature-convention}
 \int_{D_a}\dd\alpha_a=-2\pi D_a^2
\end{equation}
である．

各 $D_a$ から交点の小円板を除いた genus $g_a$，$n_a$ 境界成分の曲面を $P_a$ とする．
$P_a$ 上に，境界 collar で radial circles を葉とし，
内部には $2g_a+n_a-2$ 個の simple saddles だけを持つ実数値 Morse function の
connected-level-set 葉層 $\mathcal F_a$ を指定する．
各 connected critical leaf は一つの saddle を含む．
$n_a\geq2$ なら，このような Morse function は存在する．
例えば genus 零の cobordism に split--merge saddle pair を $g_a$ 組挿入すればよい．
$g_a=0,n_a=2$ の場合は annulus 上の正則な円周葉層である．
marking は，交点のラベルと，実際の $P_a$ 上の葉層の相対 isotopy class を含む．
その Reeb graph を $T_a$ と書く．$b_1(T_a)=g_a$ であり，有理成分の場合に限り tree である．
各交点 $p\in D_a\cap D_b$ に対応する $T_a,T_b$ の葉を同一視し，
そこでできる二価頂点を必要に応じて除いた graph を $\Gamma$ とする．
全成分が有理で二価なら，$\Gamma$ は円周である．

\begin{theorem}[SNC 因子の有限相対構成]
\label{snc:main}
仮定 \eqref{snc:valence}--\eqref{snc:positive-class} の下で，
$\omega$ と同じ cohomology class の Kähler form $\omega'$ と，
連結 fiber を持つ proper continuous surjection
\[
 \pi:(M,\omega')\longrightarrow\operatorname{Cone}_{\rm cl}(\Gamma)
\]
を次の性質を持つように構成できる．
\begin{enumerate}
\item cone point の fiber は $W=M\setminus D$ の，適切に選んだ
  logarithmic Weinstein skeleton $\Sigma$ である．
\item $W\setminus\Sigma$ の正則 fiber は Lagrangian $\TT^2$ であり，
  $\Gamma$ の三価頂点から延びる開放射辺上の fiber は
  $S^1\times\figeight$ である．他の内部 singular-fiber locus はない．
\item outer graph $\Gamma$ の逆像は正確に $D$ である．
  $D_a$ 上の connected fibers は指定した $\mathcal F_a$ の葉と一致し，
  crossing collars では radial circles，交点では一点となる．
\item 各交点の十分小さい holomorphic charts で，両 divisor を座標軸とする
  正の diagonal product Kähler germs をあらかじめ指定できる．
  collar の normal-circle action germs はそれらの product charts と一致する．
  各成分の normal connection の holonomy は構成中に選択する．
  collar の葉層は，その選択後の normal charts を用いて，固定した $D_a$ 上の
  $\mathcal F_a$ を lift したものである．
\item 各 $D_a$ の有限個の指定点が $\mathcal F_a$ の互いに異なる正則葉上にあり，
  crossing collars の外にあるなら，それらの周囲に divisor-compatible な
  holomorphic flat Kähler balls を同時に用意できる．
  これらの点は regular elliptic-boundary patches に含まれる．
\item $\Sigma$ を選んだ後の end の変更は，その近傍の Liouville structure を変えず，
  新しい trapping を作らない．全 singular fibers は，Evans--Mauri の意味で
  isotropic strata を持つ．特定の skeleton や Whitney 正則性は指定しない．
\end{enumerate}
個々の $D_a$ の ampleness，$D_a^2$ の正値性，
$H^1(M;\RR)=0$，または $Q=(D_a\cdot D_b)$ の非退化性を追加仮定する必要はない．
\end{theorem}

\begin{corollary}[多直線配置の marked な有限 compactification]
\label{comp:finite}
$M=\PP^2$，$D=L_1\cup\cdots\cup L_m$ を一般位置の $m\geq4$ 本の直線の和とし，
$c>0$ と各 $L_a$ 上の marked Morse 葉層を固定する．
ある Kähler form $\omega'\in cH$ に対して，定理\ref{snc:main}の構成が得られる．
各成分の Reeb graph は $m-1$ 個の labelled leaves と $m-3$ 個の三価頂点を持つ
tree であり，$\Gamma$ は対応する葉を交点ごとに貼り合わせた graph である．
crossing germs，full saddle collars，および有限個の指定された flat patches も
同時に保持できる．
\end{corollary}
\begin{proof}
$2\pi\sum_a w_a=c$ となる正の weights を選べば
$[\omega]=2\pi\sum_a w_a[L_a]=cH$ である．
全 $g_a=0$，$n_a=m-1\geq3$ なので，定理を適用できる．
\end{proof}

形式を $\omega'$ に替える点には実質的な意味がある．
固定した $\omega$ と固定した $D$ に対して，交点が非直交であるまま
commuting normal-circle actions を要求することはできない．
この限界と，primitive の重みを変更する場合の限界を後で明示する．

\begin{lemma}[連結性と affine 性]
\label{snc:affine}
\eqref{snc:positive-class} から $D$ は連結であり，$W$ は affine である．
\end{lemma}
\begin{proof}
$D$ が二つの非空な disjoint unions に分かれるなら，
対応する正の実 divisor を $A,B$ として $A\cdot B=0$ となる．
$[\omega]=2\pi(A+B)$ と Kähler positivity より
$A^2=[\omega]\cdot A/(2\pi)>0$，$B^2>0$ である．
これは Hodge index theorem に反する．従って $D$ は連結である．

Kähler cone は開いているので，$w$ を十分近い正の有理ベクトル $v$ に替えても
$2\pi\sum v_a[D_a]$ は Kähler class である．
分母を払うと，同じ support $D$ を持つ effective integral divisor
$A_D$ で，$\mathcal O_M(A_D)$ が正の curvature metric を持つものを得る．
従って $A_D$ は ample である．十分大きい倍を very ample とし，
対応する section の零集合を hyperplane section として実現すると，
その補集合 $W$ は affine となる．
\end{proof}

\subsection{重みを固定した global logarithmic primitive}

各 $\mathcal O_M(D_a)$ に任意の smooth Hermitian metric を取り，
canonical section を $s_a$ とする．
$-\dd\dd^c\log\|s_a\|^2$ は $W$ 上の式から $M$ 上の smooth closed
$(1,1)$-form $\Theta_a$ に延び，$[\Theta_a]=2\pi[D_a]$ である．
個々の $\Theta_a$ が正である必要はない．
compact Kähler manifold の $\dd\dd^c$-lemma より，smooth real function $u$ を用いて
\[
 \omega=\sum_a w_a\Theta_a+\dd\dd^c u
\]
と書ける\cite[Lemma (5.11), pp.~266--267; Remark (5.14)]{DGMS}．従って
\begin{equation}\label{snc:log-potential}
 \psi_0=-\sum_a w_a\log\|s_a\|^2+u,\qquad
 \lambda_0=\dd^c\psi_0,\qquad
 \dd\lambda_0=\omega|_W
\end{equation}
である．重みが全て正なので $\psi_0$ は exhaustion であり，
$\dd\dd^c\psi_0=\omega>0$ だから strictly plurisubharmonic である．
$D_a$ の正の meridian $\mathfrak m_a$ に対して
\begin{equation}\label{snc:residues}
 \lim_{\rho\to0}\int_{\mathfrak m_a(\rho)}\lambda_0=-2\pi w_a.
\end{equation}

SNC chart では $\psi_0$ は $-\sum w_a\log|z_a|^2$ と smooth function の和である．
そのため critical locus は compact であり，compact subset of $W$ 内の
小さい potential perturbation で Morse にできる．
この変更は $M$ 上の Kähler class を変えず，divisor の近傍にも影響しない．
全 critical points を含む固定 regular sublevel を compact Weinstein domain として取り，
補題\ref{ke:compact-relative-morse-smale}を適用する．これは
定理\ref{snc:main}に依存しない compact Weinstein domain の補題であり，
この時点の二形式と境界 collar の primitive を変えずに，内部に支持された
exact correction $dQ$ によって Morse--Smale にできる．この補正を外側で零に延ばし，
以後 $\lambda_0$ にはこの $dQ$ を含める．End の logarithmic primitive と全 periods
は変わらず，同補題により skeleton は次元順の閉 isotropic filtration を持つ．
その logarithmic Weinstein skeleton を $\Sigma$ とする．
$\Sigma$ は $W$ の homotopy type を持つ connected compact set であり，
この閉 filtration は Evans--Mauri\cite[Definitions~2.1 and 2.5]{EM}の意味での
isotropic strata を与える．Whitney 条件は用いない．

次に，全 critical values より高い end の中で，交点の Kähler form を product にする．
同時に，指定された divisor 上の有限個の正則点を中心とする互いに素な charts を取り，
divisor が一つの座標軸となる holomorphic flat balls を作る．
holomorphic chart の potential を $\phi=Q+O(r^3)$ と書き，
所望の正の diagonal quadratic form $Q'$ に対して
\[
 \phi'=(1-\chi(\log r))\phi+\chi(\log r)Q'
\]
を用いる．transition を長い logarithmic annulus に取ると
$h(r)=\chi(\log r)$ に対する $r|h'(r)|+r^2|h''(r)|$ を小さくでき，
$Q'-\phi=O(r^2)$ なので cutoff error は正の二形式の凸結合に対して任意に小さくなる．
従って Kähler positivity を保ち，cohomology class も変えない．
補正の微分は $O(r)$ であり，logarithmic normal differential は end で発散するので，
chart を縮めれば新しい critical point は生じない．
変更領域の potential value を元の全 critical values より高く保つため，
$\Sigma$ も変わらない．
以下ではこれまでの変更後の form と primitive を再び $\omega',\lambda_0$ と書く．
この変更は指定した $D_a$ 上の foliation を動かさない．
正則葉上の点では，normal-circle action と葉の局所 transverse coordinate が
regular elliptic-boundary model を与える．
ambient form が holomorphic flat であることと，
fibration を symplectic coordinates でその標準形にすることは区別して用いる．

crossing product charts を固定する relative symplectic tubular charts を選ぶ．
normal first derivative は complex-linear に取る．
normal coordinate の変換は $z\mapsto a z+O(|z|^2)$，$a\neq0$ であるから，
angular pole の差を引いた remainder は有界である．
交点では二つの charts が初めから product として一致する．
この regularization と以下の connection の変更は，
標準的な SNC symplectic neighborhoods の枠組みに属する
\cite[\S2.4, Appendix A]{EM}\cite{TMZ}．
ただし\cite[Theorem 2.13]{TMZ}は一般には symplectic form の変形を許す．
ここで固定した $\omega'$ に対する charts は，既に product とした crossing germs に
相対的な tubular Moser により得る．実際，normal derivative を complex-linear かつ
symplectic に選べば，tubular pullback と coupling form は zero section 上で二形式として一致する．
その差の radial homotopy primitive は二次以上で消え，固定 product charts では零に取れる．
十分薄い近傍で両形式の線分は非退化なので，Moser field は zero section と normal first jet，
および crossing charts を固定する．従ってこの段階で actual $\omega'$ をさらに変形しない．

\subsection{任意の自己交点数に対する curvature balance}

$Q_{aa}=D_a^2$，$Q_{ab}=|D_a\cap D_b|$ $(a\neq b)$ とし，
$A_a=\int_{D_a}\omega'$ とする．仮定から
\begin{equation}\label{snc:area-balance}
 A_a=2\pi\sum_b Q_{ab}w_b
 =2\pi\left(w_aD_a^2+\sum_{p\in D_a\cap D_b}w_b\right).
\end{equation}
最後の和は全ての交点 $p$ について取り，同じ $b$ が複数回現れ得る．

各 crossing cap $p\in D_a\cap D_b$ で
$\omega'|_{D_a}=\dd I\wedge\dd\theta_b$，$p=\{I=0\}$ と書く．
$P_a$ 全体には
\[
 \kappa_a=-\omega'|_{D_a}/w_a
\]
を指定する．cap では normal angular frame $\theta_a$ を用い，
\begin{equation}\label{snc:connection-cap}
 \alpha_a=\dd\theta_a+A_{a,p}(I)\dd\theta_b,\qquad
 A_{a,p}=0\text{ near }I=0,\qquad
 A_{a,p}=(w_b-I)/w_a\text{ near }P_a
\end{equation}
と指定する．この curvature の全積分は
\begin{equation}\label{snc:curvature-integral}
 \int_{D_a}\kappa_a
 =\frac{2\pi\sum_{p\in D_a\cap D_b}w_b-A_a}{w_a}
 =-2\pi D_a^2.
\end{equation}
従って，これは $N_{D_a/M}$ の connection として実現できる．
$D_a^2\leq0$ でも，cap 内の curvature が必要な分を補償するため問題はない．

$P_a$ 上で $\beta_a=0$，cap では
\[
 \beta_a=(I-w_b+w_aA_{a,p}(I))\dd\theta_b
\]
とする．$P_a$ との overlap では $\beta_a=0$ であり，
$\dd\beta_a=\omega'|_{D_a}+w_a\kappa_a$ である．
normal bundle の coupling model は
\begin{equation}\label{snc:coupling}
 \Omega_a=\omega'|_{D_a}+\dd(\mu_a\alpha_a),\qquad
 \lambda_a=(\mu_a-w_a)\alpha_a+\beta_a.
\end{equation}
$\dd\lambda_a=\Omega_a$ であり，十分薄い collar では $\Omega_a$ は symplectic である．
交点では両側の式が
\[
 (\mu_a-w_a)\dd\theta_a+(\mu_b-w_b)\dd\theta_b
\]
に一致するので，これらの primitive は一つの
$\lambda_{\rm col}$ に貼り合わさる．

ここで connection を選ぶ自由は，実際の $\omega'$ の collar を失わせない．
二つの connections の差を $\pi^*\eta_a$ と書き，
$\alpha_{a,s}=\alpha_{a,0}+s\pi^*\eta_a$ とする．
対応する coupling forms の相対 Moser field を
\[
 \iota_{Y_s}\Omega_{a,s}=-\mu_a\pi^*\eta_a
\]
で定める．crossing product charts では $\eta_a=0$ と取れる．
$Y_s$ は zero section とそこでの first jet を固定し，
normal-circle に不変で，$Y_s\mu_a=0$ である．
よって縮めた collar 上で，connection の変更を equivariant，moment-preserving な
symplectomorphism に吸収できる．以下の式はこの同定後の実際の collar 上の式とみなす．

\begin{lemma}[full Morse 葉層を保つ Liouville field]
\label{snc:collar-dynamics}
$\iota_Z\omega'=\lambda_{\rm col}$ とする．$P_a$ 上では
\[
 Z\mu_a=\mu_a-w_a,\qquad Z^{\rm horizontal}=0.
\]
cap では $A=A_{a,p}$ と略記すると
\begin{equation}\label{snc:cap-dynamics}
 Z\mu=\mu-w_a,\qquad
 ZI=\frac{I-w_b+w_aA(I)}{1+\mu A'(I)}.
\end{equation}
従って normal circles と $\mathcal F_a$ の lift からなる collar 葉層を，
saddle fiber 全体も含めて $Z$ の流れが保つ．
\end{lemma}
\begin{proof}
$P_a$ 上では
$\Omega_a=(1-\mu_a/w_a)\omega'|_{D_a}+\dd\mu_a\wedge\alpha_a$
である．よって $Z=(\mu_a-w_a)\partial_{\mu_a}$ である．
cap では
\[
 \lambda_{\rm col}=(\mu-w_a)\dd\theta_a+
     (I-w_b+\mu A(I))\dd\theta_b
\]
を \eqref{snc:coupling} に代入して \eqref{snc:cap-dynamics} を得る．
normal radius を小さくすれば $1+\mu A'(I)>0$ である．
cap の radial circle leaves と $P_a$ 上の全 connected critical leaves が保存される．
$g_a=0,n_a=2$ では $P_a$ は annulus であり，同じ計算がその正則円周葉層を保つ．
\end{proof}

\subsection{plumbing，component holonomy，graph periods}

$N$ を $D$ の小さい plumbing neighborhood，$S=\partial N$ とする．
$N$ は $D$ に retract する．曲面を交点で貼り合わせる Mayer--Vietoris 計算により
\[
 H_2(N;\ZZ)\simeq\ZZ^r,\qquad H_3(N;\ZZ)=0,
\]
かつ
\begin{equation}\label{snc:divisor-homology}
 0\longrightarrow\bigoplus_aH_1(D_a;\ZZ)
 \longrightarrow H_1(D;\ZZ)
 \longrightarrow H_1(G;\ZZ)\longrightarrow0
\end{equation}
を得る．右端は free なのでこの列は非標準的に分裂し，$H_1(N;\ZZ)$ も free である．
複数交点は $G$ の別々の辺として数える．

\begin{proposition}[meridian と divisor の完全列]
\label{snc:integer-sequence}
整数係数で
\begin{equation}\label{snc:integer-exact}
 0\longrightarrow\operatorname{coker}(Q:\ZZ^r\to\ZZ^r)
 \longrightarrow H_1(S;\ZZ)
 \longrightarrow H_1(D;\ZZ)\longrightarrow0
\end{equation}
が成り立つ．左側の subgroup は normal meridians で生成され，
その関係式は $Q$ の行で与えられる．
特に $Q$ が特異でもよく，cokernel の torsion も保持される．
有理成分だけの場合には $H_1(D;\ZZ)=H_1(G;\ZZ)$ となる．
\end{proposition}
\begin{proof}
Poincaré--Lefschetz duality と universal coefficient theorem より
$H_2(N,S;\ZZ)\simeq H^2(N;\ZZ)\simeq\ZZ^r$，
$H_1(N,S;\ZZ)\simeq H^3(N;\ZZ)=0$ である．
$H_2(N)\to H_2(N,S)$ は intersection pairing，すなわち行列 $Q$ である．
対の完全列の該当部分の cokernel を取れば \eqref{snc:integer-exact} を得る．
normal disk の境界が meridian なので，生成元と関係式の説明も従う．
\end{proof}

各 $T_a$ は $n_a$ 個の一価頂点と $2g_a+n_a-2$ 個の三価頂点を持つ．
従って $b_1(T_a)=g_a$ であり，これらを交点ごとに貼り合わせると
\begin{equation}\label{snc:graph-counts}
 b_1(\Gamma)=\sum_a g_a+e-r+1,\qquad
 |V_3(\Gamma)|=\sum_a(2g_a+n_a-2)=2\sum_a g_a+2e-2r.
\end{equation}
$D\to\Gamma$ は $H_1$ の全射を誘導するが，正の genus の成分があると
一般には同型でない．component 上の $2g_a$ 個の periods 全てを消すには，
graph から引き戻した basic forms だけでは不足し得る．
この不足は次の holonomy の自由度で解消する．

\begin{lemma}[Reeb 商で消える正則葉の cycles]
\label{snc:reeb-leaf-kernel}
Cap へ radial に延長した Reeb graph を，辺を延ばす前の graph と同じ位相型で同定し，
商写像を $q_a:D_a\to T_a$，$q:D\to\Gamma$ とする．
整数格子
\[
 L_a=\ker\bigl((q_a)_*:H_1(D_a;\ZZ)\longrightarrow H_1(T_a;\ZZ)\bigr)
\]
は rank $g_a$ の primitive Lagrangian sublattice である．
$T_a$ の maximal subtree の外の $g_a$ 本の辺から正則葉を一つずつ選ぶと，
その向き付き円周 $\ell_{a,1},\ldots,\ell_{a,g_a}$ が $L_a$ の基底になる．
また，次は完全列である：
\begin{equation}\label{snc:reeb-homology-sequence}
 0\longrightarrow\bigoplus_aL_a\longrightarrow H_1(D;\ZZ)
 \xrightarrow{q_*}H_1(\Gamma;\ZZ)\longrightarrow0.
\end{equation}
従って実係数では
\begin{equation}\label{snc:leaf-obstruction-space}
 \mathcal O_{\mathcal F}(D):=
 H^1(D;\RR)/q^*H^1(\Gamma;\RR)
 \cong\bigoplus_a\operatorname{Hom}(L_a,\RR),
 \qquad \dim\mathcal O_{\mathcal F}(D)=\sum_a g_a.
\end{equation}
$q^*H^1(\Gamma;\RR)$ の各 class は，regular annuli 内に支持され，
全 critical leaves と crossing points の近傍で零の smooth closed basic form で表せる．
\end{lemma}

\begin{proof}
閉曲面の simple Morse function に対する，level periods の消滅と smooth basic
representative の対応は既知である\cite[Lemma~4.8, p.~2421; \S5, pp.~2424--2425]{IKM}．
ここでは本節の marked pants decomposition に対し，整数格子と SNC 接続を直接確認する．
選んだ $g_a$ 本の正則円周で $D_a$ を切ると，pants と cap disks を貼る graph は tree
になるため，切断曲面は連結で genus 零である．切断後の各円周の両岸を結ぶ arcs を
戻せば，これらの円周に intersection pairing で双対な cycles が得られる．
従って円周の classes は独立で primitive であり，互いに素なので交点数は全て零である．
Rank は $g_a$ なので Lagrangian sublattice を張る．
一方，graph の各 loop は pants 内の arcs で閉曲面上へ lift できるから $(q_a)_*$ は全射である．
$b_1(T_a)=g_a$ と rank を比較すれば，選んだ格子がその kernel 全体である．

$H_1(D)$ の列\eqref{snc:divisor-homology}と，$T_a$ を交点ごとに貼る graph の列
\[
 0\longrightarrow\bigoplus_aH_1(T_a;\ZZ)
 \longrightarrow H_1(\Gamma;\ZZ)\longrightarrow H_1(G;\ZZ)\longrightarrow0
\]
を比較する．右端の写像は恒等であり，左端は全射なので，kernel の列から
\eqref{snc:reeb-homology-sequence}を得る．実係数で双対化すれば
\eqref{snc:leaf-obstruction-space}が従う．

最後に，$\Gamma$ 上の cellular $1$-cochain で class を表す．各辺の内部に短い interval
を取り，その cochain の値を積分に持つ bump form $g(u)\dd u$ を置く．
全 vertices と指定された有限個の regular values を避けて支持を選べる．
葉の商へ引き戻し，支持の外で零に延ばせば，crossings と全 critical leaves の近傍でも
滑らかな closed basic form が得られる．同じ式を normal direction に一定に延ばせば，
plumbing collar 上でも全 fibers の tangent spaces を消す closed form となる．
\end{proof}

\begin{lemma}[connection holonomy による periods の制御]
\label{snc:holonomy}
$\eta_a$ を $D_a$ 上の closed real one-form で，全 crossing disks の近傍で $0$ とする．
connection を $\alpha_{a,s}=\alpha_a+s\pi^*\eta_a$ に替え，
$\beta_a$ は変えずに $\Omega_{a,s},\lambda_{a,s}$ を \eqref{snc:coupling} で定める．
十分薄い collar 上に，crossing germs と zero section を固定する Moser isotopy
$\Phi_s$ が存在し，$\Phi_s^*\Omega_{a,s}=\Omega_{a,0}$，$\mu_a\circ\Phi_s=\mu_a$ である．
その primitive の class の変化は
\begin{equation}\label{snc:holonomy-variation}
 [\Phi_s^*\lambda_{a,s}-\lambda_{a,0}]
     =-s w_a[\pi^*\eta_a].
\end{equation}
この操作は $D_a$ 上の marked foliation を固定し，新しい normal charts 上に
その同じ foliation を lift した collar fibration を保つ．
\end{lemma}
\begin{proof}
$\eta_a$ が closed なので curvature と $\beta_a$ は変わらない．
$\iota_{Y_s}\Omega_{a,s}=-\mu_a\pi^*\eta_a$ と置くと，
$Y_s\mu_a=0$，$\alpha_{a,s}(Y_s)=0$ であり，horizontal part は $O(\mu_a)$ である．
従って十分薄い collar では時間 $0\leq s\leq1$ の flow が存在し，
zero section，そこでの normal first derivative，および crossing germs を固定する．
Cartan の公式から
\begin{align*}
 \frac{\dd}{\dd s}\Phi_s^*\lambda_{a,s}
 &=\Phi_s^*\left((\mu_a-w_a)\pi^*\eta_a
                +\iota_{Y_s}\Omega_{a,s}
                +\dd(\lambda_{a,s}(Y_s))\right)\\
 &=-w_a\Phi_s^*(\pi^*\eta_a)
      +\dd\bigl(\Phi_s^*(\lambda_{a,s}(Y_s))\bigr).
\end{align*}
$\Phi_s$ は identity と isotopic なので，cohomology を取って積分すると
\eqref{snc:holonomy-variation} を得る．
全 crossing disks で identity なので，他成分の model とそのまま貼り合わさる．
zero section 上も identity であり，$P_a$ 上では引き続き
$\kappa_a=-\omega'|_{D_a}/w_a$，$\beta_a=0$ である．
従って新しい connection の coupling charts では
$Z=(\mu_a-w_a)\partial_{\mu_a}$ であり，fixed marked foliation の lift を保つ．
この chart を $\Phi_1$ で actual collar に引き戻せばよい．
$\Omega_{a,s}$ は補助的な normal-bundle model であり，この引き戻し後の形式は
最初の $\Omega_{a,0}$ である．従って actual manifold 上の Kähler form $\omega'$ 自体は
この holonomy 補正では変更しない．
\end{proof}

$H^1(D_a;\RR)$ の任意の class は crossing disks と，
指定した flat points を中心とする有限個の小円板の近傍で $0$ の closed form で表せる．
実際，disks 上での局所 primitive を cutoff して引けばよい．
これを他成分側で $0$ として plumbing に延長すると，
\eqref{snc:divisor-homology} の component part を任意に制御できる．
大きい holonomy の補正が必要な場合でも，その補正を選んだ後に collar を縮めれば，
Moser isotopy の存在に支障はない．

$\lambda_{\rm col}$ と $\lambda_0$ は同じ $\omega'$ の primitive であり，
同じ重み $w_a$ を用いたので meridian residues が一致している．
その差の meridian integral は radius に依存せず，radius $0$ の極限では $0$ になる．
従って全 meridian periods は正確に $0$ である．
\eqref{snc:integer-exact} を実数係数で双対化すると，差の class は
$H^1(D;\RR)$ に由来する．torsion は real closed one-form の period では検出されない．
まず補題\ref{snc:holonomy}で各 $D_a$ 上の全 periods を消す．
その後に残る class は \eqref{snc:divisor-homology} により $H^1(G;\RR)$ から来る．
以下では，この holonomy 補正後の実際の collar primitive を再び
$\lambda_{\rm col}$ と書く．

$G$ の spanning tree の外の各辺に一つの regular cap arm を対応させ，
そこで支持された basic form $g(I)\dd I$ の積分を選ぶ．
これら $e-r+1$ 個の forms は残る graph periods を任意に補正できる．
その和を $\eta$ とすると
\begin{equation}\label{snc:exact-correction}
 \lambda_{\rm col}-\eta-\lambda_0=\dd f
\end{equation}
となる．$\eta$ は全 collar fibers の tangent spaces を消す．
その補正は Liouville field の angular components だけを変え，
\eqref{snc:cap-dynamics} は変えない．
各 cap arm の端では $\eta=0$ とし，normal direction に同じ式を延長するため，
$\eta$ は end 全体の smooth closed form である．

normal first derivatives を complex-linear に選んだことから，
両 primitive の angular poles を差し引いた remainder は有界である．
holonomy の Moser isotopies も normal first derivatives を固定するので，
この有界性を保つ．従って \eqref{snc:exact-correction} の $\dd f$ は有界であり，
$f$ の定数を一つ固定すれば $f$ も有界である．
後者は，finite plumbing cover 内で base point から各点まで
一様に有限な長さの経路を選び，既に全 period が消えていることから従う．
この議論で $H^1(M;\RR)$ の消滅は使わない．

\subsection{core を保つ cutoff}

十分高い level で $0$ から $1$ へ移る cutoff $\chi(\psi_0)$ を取り，
\[
 \lambda=\lambda_0+\dd\bigl(\chi(\psi_0)f\bigr)
\]
とする．$\dd\lambda=\omega'$ であり，core の近傍では $\lambda_0$，
十分 end 側では $\lambda_{\rm col}-\eta$ と一致する．
$\iota_{X_H}\omega'=-\dd H$ とすると
\begin{equation}\label{snc:cutoff-estimate}
 Z\psi_0=Z_0\psi_0-\chi(\psi_0)X_f\psi_0.
\end{equation}
$\chi'$ を含む項は $X_{\psi_0}\psi_0=0$ により消える．
normal distance を $\rho_a$ と書けば，SNC charts と $\omega'$ の正値性から
\[
 Z_0\psi_0\geq C\sum_a\rho_a^{-2},\qquad
 |X_f\psi_0|\leq C'\left(1+\sum_a\rho_a^{-1}\right)
\]
となる．よって cutoff を十分 end 側に置くと，変更領域の全てで $Z\psi_0>0$ となる．
新しい zero，trapping，または低い critical region への正時間の戻りは生じない．
したがって nonescaping set は正確に元の $\Sigma$ のままである．

adapted Liouville primitive の一般的存在は
\cite[Proposition 2.10]{BSV}にも与えられている．
ここで加えたのは，指定した full Morse 葉層を保つ curvature と holonomy の選択，
component および graph の全 periods の補正，core の germ を固定した cutoff，
および後述の finite endpoint の制御である．

\subsection{actual contact link と finite endpoint}

$P_a$ と transition caps では $\mu_a=\eps_a<w_a/2$ とし，
crossing product charts の内部だけでこれらの hypersurfaces の corners を丸める．
全ての cuts は collar fibers で飽和させる．
$Z\mu_a=\mu_a-w_a<0$ なので，得られる link $S_\infty$ は $Z$ に横断的であり，
$\alpha=\lambda|_{S_\infty}$ は contact form である．
normal-circle lift と toric orbits はそこで
\[
 S_\infty\longrightarrow\Gamma_S\simeq\Gamma
\]
という $\dd\alpha$-Lagrangian fibration を与える．
full saddle fiber は $S^1\times\figeight$ である．

$\Sigma$ 外の内側の全ての点は正時間の flow により $S_\infty$ に達する．
低い sublevel では logarithmic Weinstein structure，
高い end では \eqref{snc:cutoff-estimate} の nontrapping を使う．
従って内側の $\Sigma$ の補集合は
$((-\infty,0]\times S_\infty,\dd(e^t\alpha))$ と同定され，
$(t,x)\mapsto(e^t,f(x))$ が cone fibration を与える．
$\Sigma$ 全体を cone point に送る延長は，Evans--Mauri の cone-extension argument
\cite[Lemma 2.16]{EM}によって連続である．
これは completion の議論であり，有限体積側への延長は以下で別途確認する．

normal moment は
\begin{equation}\label{snc:time}
 \mu_a(t)=w_a+e^t(\mu_a(0)-w_a),\qquad
 T_a=\log\frac{w_a}{w_a-\mu_a(0)}
\end{equation}
を満たす．$P_a$ 上では曲線方向の点が固定されるので，
$T_a$ に normal circle を潰せば指定された葉に到達する．

transition cap $I\in[a,b]$，$a>0$ では，少し大きい interval に buffer を取り，
$0\leq\mu\leq\eps_a<w_a/2$ において $|ZI|\leq C$ と抑える．
\[
 T_a\leq\frac{2\eps_a}{w_a},\qquad
 |I(t)-I(0)|\leq\frac{2C\eps_a}{w_a}\quad(0\leq t\leq T_a).
\]
$\eps_a<a w_a/(4C)$ とし，上側の buffer からも退出しないように縮めれば，
$\mu=0$ に達するまで $I\geq a/2>0$ である．$C=0$ ではこの制約は不要である．
従って別の divisor に先に到達することはない．
全ての caps が有限個であるため，一つの十分薄い collar で同時に処理できる．

さらに $J=I-w_b+\mu A(I)$ と置くと $ZJ=J$ であり，
\begin{equation}\label{snc:endpoint}
 I_{\rm end}
 =\frac{w_aI_0-\mu_0w_b+w_a\mu_0A(I_0)}{w_a-\mu_0},\qquad
 \frac{\partial I_{\rm end}}{\partial I_0}
 =\frac{w_a(1+\mu_0A'(I_0))}{w_a-\mu_0}>0.
\end{equation}
第二式では transition cap の link 上で $\mu_0=\eps_a$ が一定であることを使う．
$P_a$ 側では $I_{\rm end}=I_0$ である．

product crossing では二つの actual normal moments が同時に存在する．
ここでのみ最初の collapse time を $\min(T_a,T_b)$ と書く．
$R=(w_a-\mu_a)/(w_b-\mu_b)$ は flow invariant であり，endpoint は
\[
 \begin{cases}
 (0,w_b-w_a/R),&R\geq w_a/w_b,\\
 (w_a-w_bR,0),&R\leq w_a/w_b
 \end{cases}
\]
となる．丸めた link arc 上で $R$ は厳密に単調なので，
その quotient は二本の divisor arms の quotient と同相である．
$R=w_a/w_b$ では両 circle が同時に潰れ，交点に到達する．

従って endpoint map $\Gamma_S\to\Gamma_D$ は，$P_a$ 上で恒等，
transition arms 上で厳密単調，crossing で同相である．
同じ flow による写像なので overlaps で一致し，graph 全体の同相を与える．
最初の到達時刻 $T(q)$ は正の連続関数である．
saddle の近傍では一定で，product crossing では二つの連続関数の最小値である．
base は $-\infty<t\leq T(q)$ に cone point を加えたものとなり，
\[
 (t,q)\longmapsto(e^{t-T(q)},q)
\]
で $\operatorname{Cone}_{\rm cl}(\Gamma)$ と同相になる．
この同相による閉端での連続性も actual collar 座標で確認できる．
$r=e^{t-T(q)}$ と置くと，component face と transition cap では
$r=1-\mu_a/w_a$，product crossing では
\[
 r=\max\{1-\mu_a/w_a,\,1-\mu_b/w_b\}
\]
である．従って $D$ への任意の接近で $r\to1$ となる．
Graph coordinate は component 上では固定葉の quotient，transition cap では
\eqref{snc:endpoint}の厳密単調な endpoint coordinate により連続に延びる．
Product crossing では不変量 $R$ を用い，交点への接近で $R\to w_a/w_b$ となるため，
同じ crossing value へ収束する．Normal charts の同定は zero section で恒等なので，
この延長は実際の $D$ 上の指定された全葉に一致する．
collapse は全て fiber ごとに行われるので，新しい fiber の切断や連結成分は生じない．
outer graph の逆像は正確に $D$ となり，boundary fibers の記述も従う．
以上で定理\ref{snc:main}を証明した．

\subsection{固定形式と固定 residue に関する限界}

\begin{proposition}[fixed form の orthogonality 障害]
\label{snc:orthogonality}
固定した $(M,\omega)$ と固定した SNC divisor $D$ に対して，
交点近傍で commuting normal-circle actions を持つ product toric collar が存在するなら，
各交点で $D_a$ と $D_b$ は $\omega$-orthogonal である．
従って定理\ref{snc:main}で $\omega'$ を任意の元の $\omega$ に置き換えることはできない．
\end{proposition}
\begin{proof}
交点の tangent space において，第一の circle の固定部分空間は $TD_a$，
非自明な weight space はその normal plane である．
$\omega$ の invariance により，この二つの weight spaces の mixed pairing は $0$ である．
product action では第二の divisor の tangent plane が第一の normal plane なので，
$TD_a$ と $TD_b$ は symplectically orthogonal である．
これは \cite[Proposition 2.5]{BSV} の局所線形代数でもある．

具体例として，$\PP^2$ の Fubini--Study form と
$D=\{x_1=0\}\cup\{x_2-x_1=0\}\cup\{x_0=0\}$ を取る．
各成分の valence は二で，正の weights による \eqref{snc:positive-class} を満たせる．
しかし $[1:0:0]$ の affine coordinates では最初の二直線の tangent spaces は
$\CC(0,1)$ と $\CC(1,1)$ であり，標準 Hermitian form に関して直交しない．
\end{proof}

元の $\omega$ 上に fibration を戻したければ，fixed-class Kähler path に対する
Moser isotopy を用いればよい．このとき compactifying divisor は $D$ と isotopic な
symplectic divisor へ移る．固定した holomorphic $D$ を同時に保存するとは限らない．
元の form が必要な holomorphic product germs を既に持つ場合には，
potential の Morse perturbation を $D$ から離して行えるため，
Moser を divisor の近傍で恒等に選ぶことができる．

\begin{proposition}[global residues と core germ]
\label{snc:residue-change}
$W=M\setminus D$ とし，closed one-form $\eta$ の residue を
$\operatorname{Res}_a(\eta)=-(2\pi)^{-1}\int_{\mathfrak m_a}\eta$ と定義する．
次は完全列である：
\begin{equation}\label{snc:global-residue-exact}
 H^1(M;\RR)\longrightarrow H^1(W;\RR)
 \xrightarrow{\operatorname{Res}}\RR^r
 \xrightarrow{(c_a)\mapsto\sum c_a[D_a]} H^2(M;\RR).
\end{equation}
従って，同じ $\omega|_W$ の primitive の重みを $w$ から $w'$ に変更できるための
必要十分条件は $\sum_a(w_a'-w_a)[D_a]=0$ である．
ただし $\Sigma\hookrightarrow W$ が homotopy equivalence である core の近傍で
primitive の germ を固定したまま変更できる重みは，$w'=w$ に限られる．
\end{proposition}
\begin{proof}
相対 cohomology の低次数において
$H^2(M,W;\RR)\simeq\RR^r$ であり，その生成元は各成分の relative Thom class である．
成分同士の交点は実 codimension 四なので，この次数に新たな生成元や関係は生じない．
対 $(M,W)$ の完全列を，Thom class と meridian integral で表示すると
\eqref{snc:global-residue-exact} を得る．
同じ二形式の primitive の差は closed one-form であるから，最初の主張が従う．
relative de Rham class と prescribed wrapping numbers の関係は
\cite[\S1.1, Proposition 2.10]{BSV}でも明示されている．

差の closed form が $\Sigma$ の近傍で $0$ なら，
$\Sigma\hookrightarrow W$ が $H_1$ の同型を誘導するため，その全 periods は $0$ である．
従って差は exact であり，全 residues も $0$ となる．
これは primitive の germ の固定に関する障害であって，
異なる weights による skeleton の集合としての同一性を一般に否定するものではない．
\end{proof}

特に，面積の線形条件 \eqref{snc:area-balance} だけで
$[\omega]=2\pi\sum w_a[D_a]$ という global 条件を代用してはいけない．
前者は collar の curvature balance であり，後者は complement 上の global exactness を保証する．
また，graph-period correction は weights を固定した後の操作である．
residues が異なる差を basic graph forms と bounded exact cutoff だけで消すことはできない．

\subsection{valence 条件と具体例}

$P_a$ の Euler 数は $2-2g_a-n_a$ である．
境界で radial な実数値 Morse function に extrema を許さず saddles だけを置くなら，
その数は $2g_a+n_a-2$ でなければならない．
$n_a\geq2$ なら上述の構成により実現でき，$n_a=2$ を除外する必要はない．
$n_a=0$ なら compactness により最大と最小が必要である．
$n_a=1$ なら唯一の境界成分上の値が一定なので，非定数の関数は少なくとも一つの
内部 extremum を持つ．従って $n_a\geq2$ は，ここで指定した
実数値 Morse function による無 extrema の collar 模型の条件として sharp である．
circle-valued foliations や elliptic strata を許す別の模型は，この主張の対象ではない．

\begin{corollary}[graph と core の数値]
\label{snc:counts}
定理\ref{snc:main}では
\[
 b_1(\Gamma)=\sum_a g_a+e-r+1,\qquad
 \#\{\text{開放射分岐辺}\}=2\sum_a g_a+2e-2r,
\]
\[
 \chi(\Sigma)=\chi(M)-2r+2\sum_a g_a+e.
\]
\end{corollary}
\begin{proof}
最初の二式は \eqref{snc:graph-counts} である．
三重交点がないので $\chi(D)=\sum_a(2-2g_a)-e$ である．
$W$ は有限 CW 型を持つ実四次元の orientable manifold なので
$\chi(W)=\chi_c(W)=\chi(M)-\chi(D)$ である．
$\Sigma$ が $W$ の homotopy type を持つことから最後の式を得る．
\end{proof}

\paragraph{$\PP^1\times\PP^1$ の toric boundary.}
二本ずつの horizontal fibers と vertical fibers の和を $D$ とする．
$r=e=4$，各成分の valence は二で，$G$ と $\Gamma$ は円周である．
各成分の自己交点は $0$ であり，個々の成分は ample ではない．
しかし両 ruling の coefficient が正なら，その class は Kähler であり，
四成分全てに正の weights を配ることができる．
$W=(\CC^*)^2$ で，core は $T^2$ の homotopy type を持つ．
この例は nonample components と valence 二を同時に含む．

\paragraph{二本の $(1,1)$ 曲線.}
$M=\PP^1\times\PP^1$ 上で二本の smooth $(1,1)$ curves を一般に取り，
二点で横断的に交わらせる．$r=e=2$，各 valence は二で，
$G$ は二頂点を二辺で結ぶ multigraph，$\Gamma$ は円周である．
任意の正の $w_1,w_2$ に対して weighted class は Kähler であり，
\[
 Q=\begin{pmatrix}2&2\\2&2\end{pmatrix},\qquad
 H_1(S;\ZZ)\simeq\ZZ^2\oplus\ZZ/2.
\]
実際，$Q$ の Smith normal form は $\operatorname{diag}(2,0)$ であり，
\eqref{snc:integer-exact} の右端 $H_1(G;\ZZ)=\ZZ$ は free なので列は非標準的に分裂する．
$\chi(W)=2$，$b_1(W)=1$ である．後者は \eqref{snc:global-residue-exact} と，
二成分の class がともに $(1,1)$ であることから従う．
$W$ は complex dimension 二の Stein manifold なので二次元 CW 型を持ち，
$\chi(W)=1-b_1(W)+b_2(W)$ より $b_2(W)=2$ となる．
この core を simple focus--focus fiber と同定することはできない．
少なくともその Euler 数は，simple nodal fiber の $1$ と異なる．

\paragraph{三本の $(1,1)$ 曲線.}
同じ曲面上に三本の smooth $(1,1)$ curves を一般に取り，三重交点を避ける．
各 pair は二点で横断的に交わるため $r=3$，$e=6$，全 valence は四である．
従って $b_1(\Gamma)=4$，trivalent vertices は六個，core の Euler 数は四である．
これは実際に branching を持ち，直線配置の単純グラフには含まれない例である．
正の weights の任意の選択に対して weighted class は Kähler である．

\paragraph{負の自己交点を含む例.}
Hirzebruch surface $\mathbb F_n$ の toric boundary を
$S_0+F_0+S_\infty+F_\infty$ とする．
$S_0^2=-n$，$S_\infty=S_0+nF$，$F_i^2=0$ であり，全 valence は二である．
weights を $u_0,u_\infty,v_0,v_\infty>0$ とすると，weighted divisor class は
\[
 (u_0+u_\infty)S_0+(nu_\infty+v_0+v_\infty)F.
\]
$v_0+v_\infty>nu_0$ ならこの class は ample であり，定理が適用される．
cap の curvature balance は負の $S_0^2$ を含めて
\eqref{snc:curvature-integral} により同時に解かれる．


\paragraph{二本の smooth plane cubics.}
$M=\PP^2$ で，smooth cubics $D_1,D_2$ を九点で横断的に交わるように取る．
$g_1=g_2=1$，$r=2$，$e=9$，$n_1=n_2=9$ であり，
任意の正の weights が Kähler class を与える．
各 $P_a$ の Morse foliation は九個の saddles を持ち，
\[
 b_1(\Gamma)=10,\qquad |V_3(\Gamma)|=18,\qquad \chi(\Sigma)=12.
\]
meridians の整数表示は
\[
 H_1(W;\ZZ)=\langle \mathfrak m_1,\mathfrak m_2\mid
          3\mathfrak m_1+3\mathfrak m_2=0\rangle
     \simeq\ZZ\oplus\ZZ/3.
\]
これは pair の homology exact sequence と，line が各 cubic と三点で交わることから従う．
この例では $H^1(D;\RR)$ の dimension は十二であるのに対して
$H^1(\Gamma;\RR)$ の dimension は十であり，
foliation を固定した basic forms だけでは任意の period 差を補正できないことが分かる．
補題\ref{snc:holonomy}は，二つの elliptic components の holonomy を選んでこれを解消する．

\subsection{固定 connection と全 leaf-period 条件}

以下で固定するのは unitary holonomy の値だけでなく，normal connections と
actual collar への coupling charts，およびその full fiber partition である．
許す補正は closed foliation-basic form と，core を保つ exact cutoff である．
この範囲で，先の period 補正の必要十分条件を記述できる．

\begin{proposition}[固定 coupling charts に対する相対 primitive 接着]
\label{snc:fixed-holonomy}
本節の logarithmic models の下で，$\lambda_0$ と $\lambda_{\rm col}$ は
同じ二形式 $\omega'$ の primitive であり，同じ residues $-2\pi w_a$ を持つとする．
同じ angular poles を引いた remainders が有界であることも，この model の条件に含める．
Plumbing retraction の境界への制限を $\rho:S\to D$ と書くと，
命題\ref{snc:integer-sequence}により一意な $c\in H^1(D;\RR)$ があり，
$[\lambda_{\rm col}-\lambda_0]=\rho^*c$ となる．次は同値である．
\begin{enumerate}
\item Collar の closed basic form $\eta$ と有界関数 $f$ により
  $\lambda_{\rm col}-\eta-\lambda_0=\dd f$ とできる．
  $\eta$ は全 critical leaves と crossing germs の近傍で零に選べる．
\item $c\in q^*H^1(\Gamma;\RR)$，すなわち $[c]=0$ in $\mathcal O_{\mathcal F}(D)$．
\item 補題\ref{snc:reeb-leaf-kernel}の各 $\ell_{a,k}$ に対して
  \begin{equation}\label{snc:fixed-leaf-period-test}
   \int_{\widetilde\ell_{a,k}}(\lambda_{\rm col}-\lambda_0)=0
   \qquad(1\le k\le g_a)
  \end{equation}
  が成り立つ．ここで $\widetilde\ell_{a,k}$ は collar 上の任意の circle-bundle lift である．
\end{enumerate}
必要なら固定 charts を十分薄い共通 subcollar に制限する．これらの条件の下で，
core の primitive germ，その固定 collar の全葉，normal moments，
full saddle collars と crossing germs を保つ本節の exact-cutoff 構成を実行できる．
Finite endpoint で得られる $D$ 上の全葉も元の指定と一致する．
\end{proposition}

\begin{proof}
Meridian periods が零なので，\eqref{snc:fixed-leaf-period-test}は lift の選び方に依存せず，
$c([\ell_{a,k}])$ に等しい．補題\ref{snc:reeb-leaf-kernel}より (2) と (3) は同値である．
各 $\widetilde\ell_{a,k}$ は一つの full torus fiber 内に取れるので，
basic $\eta$ と $\dd f$ の積分は零である．従って (1) なら (3) である．
逆に (2) の下では同補題の edge-bump representatives を collar に延ばして $\eta$ を取る．
差は exact となり，同じ angular poles を引いた remainder の有界性と
有限 plumbing cover の経路による議論から，$\dd f$ と $f$ を有界に取れる．
従って (1) が成立する．

Core へ接ぐ cutoff の nontrapping 評価は\eqref{snc:cutoff-estimate}と同じである．
ここでは一般の basic $\eta$ を $P_a$ の regular annuli にも置くので，
curve 上の点を固定することと，全葉を保つことを区別する．
$\lambda_{\rm col}-\eta$ の Liouville field を $Z$ とすると
\[
 Z\mu_a=\mu_a-w_a,\qquad \alpha_a(Z)=0,\qquad
 (1-\mu_a/w_a)\iota_{Z^{\rm horizontal}}\omega'|_{D_a}=-\eta.
\]
$\eta$ は葉の tangent を消すから，horizontal part は葉の tangent である．
さらに $\mu_a\le\eps_a<w_a/2$ では滑らかに有界である．
従って有限時刻\eqref{snc:time}までの base flow は同じ全正則葉上の diffeomorphism を与え，
normal circle を潰す endpoint も同じ全葉である．
Critical collars では $\eta=0$ なので full saddle leaves をそのまま保つ．
Cap の basic form $g(I)\dd I$ は angular components だけを変えるため，
\eqref{snc:cap-dynamics}と\eqref{snc:endpoint}も変わらない．
Crossing germs も支持の外であり，既存の閉端への連続性が使える．
\end{proof}

\begin{corollary}[必要な holonomy 調整の独立方向]
\label{snc:minimal-holonomy-controls}
固定された葉層に対し，$\sum_a g_a$ 本の closed real one-forms を用いる connection
変更と basic correction で，任意の $c\in H^1(D;\RR)$ の障害を消せる．
Holonomy 変更から $\mathcal O_{\mathcal F}(D)$ への線形写像の rank は正確に $\sum_a g_a$ である．
従って全ての可能な障害を制御する線形 family には，これだけの独立方向が必要十分である．
ここで holonomy 変更後の lift は新しい coupling charts を用いる．
\end{corollary}

\begin{proof}
各 $a$ で，$\int_{\ell_{a,j}}\eta_{a,k}=\delta_{jk}$ となる closed forms を選ぶ．
補題\ref{snc:reeb-leaf-kernel}の primitive lattice とその双対 cycles により選択でき，
局所 primitive の cutoff により全 crossing disks と指定 flat patches の近傍で零にできる．
$p_{a,k}=c([\ell_{a,k}])$ とし，connection を
\[
 \alpha_a\longmapsto\alpha_a+
       \sum_{k=1}^{g_a}(p_{a,k}/w_a)\pi^*\eta_{a,k}
\]
と変える．補題\ref{snc:holonomy}により障害の各座標は
$p_{a,k}-w_a(p_{a,k}/w_a)=0$ となる．残りは $q^*H^1(\Gamma)$ の class なので，
命題\ref{snc:fixed-holonomy}の basic correction で処理できる．
選んだ $\sum_a g_a$ 方向と障害座標に関する行列は，重み $-w_a$ の非零対角行列であり，
rank の主張も従う．
\end{proof}

\begin{remark}[Connection 補正の支持の制約]\label{snc:saddle-support-obstruction}
Closed form $\eta_a$ が全ての full critical-leaf collars と crossing collars の近傍で零なら，
$[\eta_a]|_{L_a}=0$ である．実際，regular annulus 上の葉の period は Stokes の定理により一定で，
端の critical collar または cap で零になる．従って，このような支持の connection 変更は
$\mathcal O_{\mathcal F}(D)$ を変えない．非零障害を消す際には，一般に全ての full saddle collars を
同時に固定することはできない．Saddle 点だけの小さい disks を固定する条件はこれより弱い．
Holonomy 補正後も $D$ 上の指定 full leaves は保たれるが，actual collar での lift は新しい
coupling charts による．障害消去後の basic graph correction は，これとは異なり
全 critical-leaf collars を支持外にできる．
\end{remark}

二本の smooth plane cubics の例では，$\dim H^1(D)=12$，$b_1(\Gamma)=10$ であり，
固定 coupling charts の障害は二つの elliptic components 上の正則葉の period 一つずつである．
各成分の二つの periods 全てを holonomy で消すことは十分だが，必要な独立調整は計二つで足りる．
全成分が有理なら $\mathcal O_{\mathcal F}(D)=0$ であり，この段階には holonomy 変更を要しない．
一方，命題\ref{ke:component-flux-space}の点固定 collar germ の障害空間とは異なる quotient
であり，その消滅を意味しない．また，任意の fibration の存在に対する必要条件や，
全ての foliation-preserving Liouville fields の分類を主張しているわけではない．

Unitary holonomy の位相値だけを固定する場合にも注意が必要である．
$\eta_a$ の全 periods が $2\pi\ZZ$ に入る変更は holonomy phases を変えないが，
\eqref{snc:holonomy-variation}の real primitive class は $-w_a[\eta_a]$ だけ動く．
この変更は normal gauge を変えるので，上の固定 coupling charts の条件には含めない．
位相値だけを固定する問題では，この整数 period の自由を別途扱う必要がある．

\subsection{実際の成分へ移す際に保持されるデータ}
\label{snc:output-data}

定理\ref{snc:main}は固定した Kähler class 内の一つの model として，
次の相対データを与える．
\begin{enumerate}
\item crossing で固定した holomorphic product germs と，
それに一致する commuting normal-circle action germs．
\item 各 labelled $D_a$ 上の実際の marked Morse 葉層と，
その connected critical leaf 全体を含む full saddle collar．
正の genus の場合の holonomy はこの構成の一部として選ぶ．
\item core から離れた有限個の divisor-compatible holomorphic flat balls と，
regular elliptic-boundary patches．局所 nodal surgery を用いる場合には，
その surgery に必要な symplectic chart との同定を別途行う．
\item core と局所 surgery の compact support を避けてさらに縮められる divisor collars．
normal moments は各 zero section 上で $0$ と正規化される．
\item core 近傍の primitive germ を固定した component および graph の period 補正．
この補正だけから，別個に選ばれた symplectic pieces の action data の一致を
結論するものではない．
\end{enumerate}
actual component への移送では，固定 crossing germs を持つ Kähler path，
divisor-relative Moser，normal-circle action germs の同定，
および compact saturated truncation 全体の transport を用いる．
これらは本節の有限 compactification の後で確認すべき別の相対条件である．
正則 band で先指定した第二作用まで合わせる条件は命題\ref{neck:regular-period-matching}，
同じ actual form 上で全葉を選んで接続する構成は第\ref{sec:gluing}節に記す．
作用適合 atlas を得るには，さらに critical action の有理直線性
（命題\ref{atl:critical-action-obstruction}）を満たすように葉族を選ぶ必要がある．
例\ref{atl:fixed-form-bent-branch}は，normal action と両端 germs の固定だけから
その直線性が従わないことを示す．
